% Part:sets-functions-relations
% Chapter: size-of-sets
% Section: enumerations-alt

\documentclass[../../../include/open-logic-section]{subfiles}

\begin{document}

\olfileid{sfr}{siz}{enm-alt}

\olsection{પરિગણનાઓ અને \usetoken{S}{enumerable} ગણો}

\begin{editorial}
  આ વિભાગ ગણસિદ્ધાંતમાં થતી રીતે પરિગણનાઓને $\Nat$ના
  (આરંભખંડો) સાથેના એક-એક વ્યાપ્ત વિધેયો તરીકે વ્યાખ્યાયિત કરે છે.
  તેથી તે \olref[enm]{sec}ની વ્યાખ્યાઓથી સહેજ અસંગત છે અને ત્યાંનાં
  બધાં ઉદાહરણોનું પુનરાવર્તન કરે છે. તે વિભાગ કરતાં આ વિભાગ થોડો
  વધુ સંક્ષિપ્ત પણ છે.
\end{editorial}

સાન્ત ગણના !!{element}sની પરિગણના કરીને આપણે તે ગણને સહેલાઈથી
નિર્દિષ્ટ કરી શકીએ. દાખલા તરીકે, ગણને આ રીતે વ્યાખ્યાયિત કરીએ છીએ:
\[
  A = \{a_1, a_2, \ldots, a_n\}.
\]
!!{element}s $a_1$, \dots, $a_n$ બધા જુદા છે તેમ ધારીએ, તો આ
$A$ અને પ્રથમ $n$ પ્રાકૃતિક સંખ્યાઓ $0$, \dots, $n-1$ વચ્ચે
!!a{bijection} આપે છે. ઊલટું, દરેક સાન્ત ગણમાં સાન્ત સંખ્યામાં જ
!!{element}s હોવાથી દરેક સાન્ત ગણને આવી સંગતિમાં મૂકી શકાય છે.
બીજા શબ્દોમાં, જો $A$ સાન્ત હોય, તો $A$ અને $\{0, \dots, n-1\}$
વચ્ચે !!a{bijection} હોય છે, જ્યાં $n$ એ $A$ના !!{element}sની સંખ્યા છે.

ચોક્કસ પ્રકારના અનંત ગણોને સ્વીકારીએ, તો કેટલાક અનંત ગણોની પણ
પરિગણના થઈ શકે તેમ સ્વીકારીશું. અનંત ગણ અને સમગ્ર~$\Nat$ વચ્ચે
!!a{bijection} હોય ત્યારે તે અનંત ગણ પરિગણિત થાય છે એમ કહીને આ
વાત ચોક્કસ કરી શકાય છે.

\begin{defn}[પરિગણના, ગણસિદ્ધાંત મુજબ]
ગણ $A$ની \emph{પરિગણના} એ એવું !!a{bijection} છે જેનો વિસ્તાર
$A$ હોય અને જેનો પ્રદેશ કાં તો પ્રાકૃતિક સંખ્યાઓનો કોઈ આરંભખંડ
$\{0, 1, \ldots, n\}$ {અથવા} પ્રાકૃતિક સંખ્યાઓનો સમગ્ર ગણ~$\Nat$ હોય.
\end{defn}

\begin{explain}
\emph{પરિગણના} શબ્દના આ ઉપયોગને સહજ આધાર છે. આપણે ગણ $A$ની
પરિગણના કરી છે તેમ કહેવાનો અર્થ એ છે કે એવું !!a{bijection} $f$
છે જેનાથી ગણ $A$ના ઘટકો ગણી શકાય છે. $0$મો ઘટક $f(0)$ છે,
1લો $f(1)$ છે, \ldots અને $n$મો $f(n)$ છે\ldots.
\footnote{હા, આપણે $0$થી ગણીએ છીએ. અલબત્ત, $1$થી પણ શરૂ કરી
શકીએ. તેનાથી કોઈ મોટો ફરક ન પડે. માત્ર~$\Nat$ને~$\PosInt$થી
બદલવું પડે.} નીચેનું ઉમેરવાથી આનું કારણ વધુ સ્પષ્ટ થઈ શકે છે:
\end{explain}

\begin{defn}
  \ollabel{defn:enumerable}
  ગણ~$A$ એ !!{enumerable} હોય જો અને માત્ર જો કાં તો
  $A = \emptyset$ અથવા $A$ની કોઈ પરિગણના હોય. ગણ~$A$ને
  !!{nonenumerable} કહીએ જો અને માત્ર જો $A$ એ !!{enumerable} ન હોય.
\end{defn}

\begin{explain}
આમ, ગણ !!{enumerable} હોય જો અને માત્ર જો તે ખાલી હોય અથવા તેના
!!{element}sને પરિગણનાથી ગણી શકાય.
\end{explain}

\begin{ex}
પ્રાકૃતિક સંખ્યાઓની પરિગણના કરતું વિધેય
$\Id{\Nat} \colon \Nat \to \Nat$ એવું તદેવ વિધેય છે, જ્યાં
$\Id{\Nat}(n) = n$. \emph{ધન} પ્રાકૃતિક સંખ્યાઓ
$\Nat^+ = \Nat \setminus \{0\}$ની પરિગણના કરતું વિધેય
$g(n) = n + 1$ છે, એટલે કે અનુગામી વિધેય.
\end{ex}

\begin{prob}
બતાવો કે ગણ $A$ એ !!{enumerable} હોય જો અને માત્ર જો કાં તો
$A = \emptyset$ અથવા કોઈ !!a{surjection} $f\colon \Nat \to A$ હોય.
બતાવો કે $A$ એ !!{enumerable} હોય જો અને માત્ર જો કોઈ
!!a{injection} $g\colon A \to \Nat$ હોય.
\end{prob}

\begin{ex}
વિધેયો $f\colon \Nat \to \Nat$ અને $g \colon \Nat \to \Nat$ને
અનુક્રમે
\begin{align*}
  f(n) & = 2n \text{ અને}\\
  g(n) & = 2n+1
\end{align*}
દ્વારા આપીએ, તો તેઓ યુગ્મ પ્રાકૃતિક સંખ્યાઓ અને અયુગ્મ પ્રાકૃતિક
સંખ્યાઓની પરિગણના કરે છે. પરંતુ એકેય !!{surjective} નથી, તેથી
એકેય $\Nat$ની પરિગણના નથી.
\end{ex}

\begin{prob}
વર્ગસંખ્યાઓ $1$, $4$, $9$, $16$, \dots{}ની પરિગણના વ્યાખ્યાયિત કરો.
\end{prob}

\begin{ex}
$\lceil x \rceil$ એ \emph{ઊર્ધ્વ પૂર્ણાંક} વિધેય હોય, જે $x$ને
ઉપરની નજીકની પૂર્ણાંક કિંમત સુધી લઈ જાય. ત્યારે આ રીતે આપેલું
વિધેય $f \colon \Nat \to \Int$:
\[
  f(n) = (-1)^{n} \left\lceil\tfrac{n}{2}\right\rceil
\]
પૂર્ણાંકોના ગણ~$\Int$ની પરિગણના નીચે મુજબ કરે છે:
\[
\begin{array}{c c c c c c c c}
f(0) & f(1) & f(2) & f(3) & f(4) & f(5) & f(6) & \dots \\ \\
\big\lceil \tfrac{0}{2} \big\rceil & -\big\lceil \tfrac{1}{2}\big\rceil &  \big\lceil \tfrac{2}{2} \big\rceil & -\big\lceil \tfrac{3}{2} \big\rceil & \big\lceil \tfrac{4}{2} \big\rceil  & -\big\lceil \tfrac{5}{2}\big\rceil & \big\lceil \tfrac{6}{2} \big\rceil & \dots \\ \\
0 & -1 & 1 & -2 & 2 & -3 & 3& \dots
\end{array}
\]
જુઓ કે $f$ ધન અને ઋણ પૂર્ણાંકો વચ્ચે આગળપાછળ “કૂદીને” $\Int$ની
કિંમતો કેવી રીતે ઉત્પન્ન કરે છે. નીચે મુજબ કિસ્સા પાડીને $f$
વ્યાખ્યાયિત થયું છે એમ પણ વિચારી શકો:
\[
f(n) = \begin{cases}
  \frac{n}{2} & \text{જો $n$ યુગ્મ હોય}\\
  -\frac{n+1}{2} & \text{જો $n$ અયુગ્મ હોય}
  \end{cases}
\]
\end{ex}

\begin{prob}
બતાવો કે જો $A$ અને $B$ એ !!{enumerable} હોય, તો $A \cup B$ પણ છે.
\end{prob}

\begin{prob}
$n$ પર ગાણિતિક અનુમાનથી બતાવો કે જો $A_1$, $A_2$, \dots, $A_n$
બધા !!{enumerable} હોય, તો $A_1 \cup \dots \cup A_n$ પણ છે.
\end{prob}

\end{document}
